\documentclass[11pt, draft]{article}
\topmargin=0mm
\oddsidemargin=0mm
\textwidth170mm
\textheight230mm
\begin{document}
%\pagestyle{empty}
\language=0
%\hyphenation{}

\begin{center}
{\bf
MORPHISMS AND BINOMIAL IDEALS IN POLYNOMIAL AND DIFFERENTIAL
POLYNOMIAL RINGS}\\
\vspace {0.3cm}
{Giuseppa Carra' Ferro}\\
{\it \small Department of Mathematics and Computer Science,
University of Catania, Viale Andrea Doria 6

95125 Catania, Italy, e-mail: carra@dmi.unict.it
}
\end{center}
\newtheorem{definition}{DEFINITION}
\newtheorem{theorem}{THEOREM}
\newtheorem{example}{EXAMPLE}
\newtheorem{remark}{REMARK}
\newtheorem{proposition}{PROPOSITION}
\newtheorem{corollary}{COROLLARY}
\vspace {1.0cm}
Binomial ideals in polynomial rings is one of the most studied subjects
in the last years, since they are simpler than generic ideals but carry
more informations than monomial ideals.
Let $I$ be a nonzero ideal in $A$=$K[X_1,\ldots,X_n]$ and let $\sigma$ be a
term ordering on $T_{A}$. Let
$F$=$\{f_{1},\ldots,f_{r}\}$ be the reduced Gr\"{o}bner basis of $I$
with respect to  $\sigma$ and let $f_{i}=\sum_{j=1,\ldots,s(i)}c_{ij}t_{ij}+
c_{i0}$, $c_{ij} \in K^{*}$, $c_{i0} \in K$, $t_{ij} \in T_{A}$ and
$t_{ij} \neq 1$ for all $i=1,\ldots,r$ and $j=1,\ldots,r(i)$.
Let $m=\sum_{i=1,\ldots,r}s(i)$. Let $\varphi_{I}$ be the map from
$B$=$K[Y_{11},\ldots,Y_{1s(1)},\ldots,Y_{r1},\ldots,Y_{rs(r)}]$=
$K[Y_1,\ldots,Y_m]$ to $A$ defined by $\varphi(Y_{ij})=t_{ij}$ for all
$i=1,\ldots,r$ and $j=1,\ldots,r(i)$.
$\varphi_{I}$ defines a ring homorphism.
Let $J_{\varphi_{I}}$=($Yij-tij$: $i=1,\ldots,r$ and $j=1,\ldots,r(i)$) in
$K[Y_{11},\ldots,Y_{1s(1)},\ldots,Y_{r1},\ldots,Y_{rs(r)},X_1,\ldots,X_n]$.
Let $\tau$ be an elimination term ordering with $X_{h} >_{\tau} Y_{ij}$ for
all $h=1,\ldots,n$, $i=1,\ldots,r$ and $j=1,\ldots,r(i))$, such that its restriction to
$T_{A}$ is $\sigma$. Let $H(I)=J_{\varphi_{I}} \cap B$.
It is well known that $H(I)$ is a prime binomial ideal and
$H(I)=Ker(\varphi_{I})$.\\
$\varphi_{I}$ defines a polynomial morphism $\phi_{I}$ from $A^{n}_{K}$
to $A^{m}_{K}$, such that
$$\phi_{I} (X_1,\ldots,X_n)=
(Y_{11},\ldots,Y_{1s(1)},\ldots,Y_{r1},\ldots,Y_{rs(r)}).$$
The closure of $\phi_{I} (A^{n}_{K})$ in the Zariski
topology is a binomial variety, whenever $\phi_{I} (A^{n}_{K}) \neq A^{m}_{K}$.
The polynomial morphism $\phi_{I}$ defines a birational isomorphism between the
affine algebraic set $V(I)$ in $A^{n}_{K}$ and an affine algebraic set
in $A^{m}_{K}$, which is intersection of a binomial affine algebraic set
and linear spaces. This result can be extended to the case of differential
affine algebraic sets.
\end{document}
